\chapter{Introduction}
\label{chap:introduction}

Writing a long document is an iterative activity. An author changes a
sentence, rebuilds the document, and looks at the result. The time between
the edit and the updated page shapes how the work feels: a short delay keeps
attention on the prose, and a long one invites distraction.

\TeX{} was designed as a batch program \citep{knuth1984texbook}, and
\LaTeX{} inherits that model \citep{lamport1994latex}. Each run reads the
whole source from the start, even when a single word has changed.

\section{Motivation}
\label{sec:motivation}

Build tools for software solved a similar problem decades ago. \emph{Make}
rebuilds only the targets whose inputs changed \citep{feldman1979make}. This
thesis asks how far the same idea carries over to typesetting, where the
output of one page can depend on every page before it.

\section{Contributions}
\label{sec:contributions}

The thesis makes three contributions:
\begin{itemize}
  \item a model of the dependencies between source files and typeset pages,
        described in Chapter~\ref{chap:design};
  \item a prototype that reuses engine state across runs; and
  \item an evaluation of rebuild latency, reported in
        Chapter~\ref{chap:evaluation}.
\end{itemize}

\section{Outline}
\label{sec:outline}

Chapter~\ref{chap:background} reviews how \TeX{} processes a document.
Chapter~\ref{chap:design} presents the design, Chapter~\ref{chap:evaluation}
evaluates it, and Chapter~\ref{chap:conclusion} concludes.
